\documentclass[10pt,a4paper]{amsart}
\usepackage[margin=19mm]{geometry}
\usepackage{amsmath,amssymb,mathtools}
\usepackage[scr=rsfs,cal=euler]{mathalfa}
\usepackage{xfrac}
\usepackage[unicode,hidelinks,bookmarks=false]{hyperref}
\newcommand{\ie}{i.e.\ }
\newcommand{\cf}{cf.\ }
\newcommand{\resp}{respectively\ }
\newcommand{\Subsection}{\subsection}
\providecommand{\nobreakdash}{\nobreak\hbox{-}}
\setlength{\emergencystretch}{2em}
\multlinegap0pt
\usepackage{bm}
\usepackage{dsfont}
\usepackage[shortlabels]{enumitem}
\usepackage{tikz}
\usepackage{stackengine}
\usepackage[normalem]{ulem}
\usepackage{algorithm}\usepackage{algpseudocode}
\newtheorem{proposition}{Proposition}
\newtheorem{theorem}{Theorem}
\newcommand{\EE}{\mathbb E}
\title{Right-Sizing Communication and Recommendation Set Size in AI-Assisted Search}
\author{Jing Dong, Prakirt Raj Jhunjhunwala, Yash Kanoria}
\date{Source: arXiv:2605.23944v1 (CC BY 4.0)}
\begin{document}
\maketitle

\noindent\emph{Source and licence.} Right-Sizing Communication and Recommendation Set Size in AI-Assisted Search. Version arXiv:2605.23944v1 (CC BY 4.0), distributed by arXiv under the Creative Commons Attribution 4.0 International licence, http://creativecommons.org/licenses/by/4.0/, as recorded in the arXiv licence metadata for this exact version and shown on the abstract page https://arxiv.org/abs/2605.23944v1. Full licence notice: \texttt{licenses/arxiv-2605.23944-CC-BY-4.0.txt}.\par\smallskip

\noindent\textit{Editorial scope.} Counted unit: the article's proposition utility\_approx (source0.tex lines 338--354). The article's complete original proof of it is delivered with it (source0.tex lines 1130--1277, one \texttt{proof} environment of 148 lines, which the article places away from the statement). No mathematics is written, repaired or re-derived: the statement and the proof are the article's own lines, in the article's own notation, and no retained line is substituted or re-flowed.  Retained uncounted as the source-local context the proof needs: lines 280--283; lines 357--364; lines 369--372; lines 487--498.  Nothing was omitted from any retained span, so the package carries no omission record.  No retained line contains a citation, so the wrapper carries no bibliography and no bibliography entry.  No retained line contains a figure environment, an image-inclusion command or drawing code, so no image is delivered and the package declares no asset dependency.  Editorial formatting only, in addition: an AMS \texttt{amsart} preamble that restates the article's own declarations and macros used by the retained lines, namely the environments \texttt{proposition}, \texttt{theorem} on separate counters, together with the standard packages those lines need.  The retained environments print in this document as Proposition 1, Theorem 1 in order of appearance, whereas the article prints the same items with the numbering of its own full document.  The complete native source of the article as distributed in the arXiv e-print for this exact version is bundled unchanged as \texttt{sources/201248/source0.tex}.
\begin{align}
\label{eq:util_representation}
    \boldsymbol\theta_i = W_i \mathbf{m} + \sqrt{1-W_i^2} \mathbf{Y}_i, \ \text{ and } \ \langle \mathbf{h} , \boldsymbol{\theta}_i \rangle = W W_i + \sqrt{1-W^2} \sqrt{1-W_i^2} X_i,
\end{align}

\begin{proposition}[Communication cost]
\label{prop: kl_approx}
Suppose $\rho$ is fixed and let $\kappa = \frac{\rho}{1-\rho^2}(d-3)$. We have,
\begin{align*}
\EE_{\mathbf{h},\mathbf{m}} \big[D_{\mathrm{KL}}\big(q_{\kappa}(\cdot|\mathbf{m}) \| p(\cdot)\big) \big] = \frac{d-2}{2} \log \frac{1}{1-\rho^2} + \mathcal{E}_{\ref{prop: kl_approx}},
\end{align*}
where $|\mathcal{E}_{\ref{prop: kl_approx}}(\rho)| \leq K_{\ref{prop: kl_approx}} \sqrt{\frac{d}{(1-\rho^2)^3}}$, and $K_{\ref{prop: kl_approx}}$ is a universal constant.
\end{proposition}

\begin{align}
\label{eq:joint_opt}
\text{OPT}_{\text{Joint}} := \max_{\rho \in [0,1), \alpha \geq 0} \ \ \Big \{ f(\rho,\alpha) - c_s \alpha + \frac{1}{2} c_c \log (1-\rho^2) \Big\},
\end{align}

\begin{theorem}[Optimality gap of sampling from a tilted distribution]
\label{thm:tilt_opt}
For a recommendation set of size $n$, the expected utility of a sampling scheme $q \in \mathcal{Q}$ is given by
\begin{align*}
U_n(q) := \mathbb{E}_{\boldsymbol\theta_1, \dots, \boldsymbol\theta_n \sim q(\cdot | \mathbf{m}), \text{ i.i.d.}} \left[ \max_{i \in [n]} \langle \mathbf{h}, \boldsymbol\theta_i \rangle \right].
\end{align*}
Then, there exists a universal constant $K_{\ref{thm:tilt_opt}}$, such that, for $d \geq 4$, the maximum expected utility achievable by any distribution in the class $\mathcal{Q}$ satisfies
\begin{align*}
\max_{q \in \mathcal{Q}} \mathbb{E}_{\mathbf{h}, \mathbf{m}} [U_n(q)] \leq \max_{q \in \mathcal{Q}^{\text{Tilt}}} \mathbb{E}_{\mathbf{h}, \mathbf{m}} [U_n(q)] + \frac{K_{\ref{thm:tilt_opt}}}{\sqrt{d-3} (1-\rho^2)},
\end{align*}
where the expectation is taken over the joint distribution of preferences $\mathbf{h}$, and messages $\mathbf{m}\sim p_{\kappa}(\cdot|\mathbf{h})$, where $\rho \in [0,1)$ and $\kappa = \frac{\rho}{(1-\rho^2)} (d-3)$.
\end{theorem}

\begin{proposition}[Product utility]
\label{prop: utility_approx}
Suppose $\rho \in [0,1)$ and $\alpha \geq 0$ are fixed and define:
\begin{align*}
    f(\rho,\alpha) := \max_{w, x \in (-1,1)} & \ \ \ \rho w + \sqrt{1-\rho^2} \sqrt{1-w^2} x \ \text{ such that } I_{\rho}(w,x) \leq \alpha,
\end{align*}
where $I_{\rho}(w,x)$ is the large deviations rate function, given by
\begin{align}
\label{eq:rate_function}
    I_{\rho}(w,x) = -\frac{\rho(w-\rho)}{1-\rho^2} - \frac{1}{2} \log(1-w^2)-\frac{1}{2} \log(1-x^2) + \frac{1}{2} \log(1-\rho^2).
\end{align}
Suppose the message precision $\kappa = \frac{\rho}{1-\rho^2}(d-3)$ and the recommendation set size $n = \lfloor e^{d\alpha}\rfloor$, then, the corresponding expected utility received by the user satisfies
\begin{align*}
    \mathbb E \left[\max_i \Big\{  W W_i + \sqrt{1-W^2} \sqrt{1-W_i^2} X_i \Big\}\right] = f(\rho,\alpha) + \mathcal{E}_{\ref{prop: utility_approx}}, \ \text{ where }  \ |\mathcal{E}_{\ref{prop: utility_approx}}| \leq K_{\ref{prop: utility_approx}} \sqrt{\frac{\log d}{d}},
\end{align*}
where $W, W_i$ and $X_i$'s are as in Eq. \eqref{eq:util_representation}, and the constant $K_{\ref{prop: utility_approx}}$ depends on $\rho$ and $\alpha$.
\end{proposition}

\begin{proof}
We prove each part of the theorem sequentially.

\paragraph{Part 1: Convergence}
We analyze the asymptotic behavior of each term as $d \to \infty$ under the scaling assumptions $\lambda_s = c_s/d + o(d^{-1})$ and $\lambda_c = c_c/d + o(d^{-1})$. Let $n = \lfloor e^{\alpha d} \rfloor$ and $\kappa = \frac{\rho}{1-\rho^2}(d-3)$ for fixed $\alpha, \rho$. From Proposition \ref{prop: utility_approx}, the expected utility of the best recommendation converges as
\begin{equation*}
   \lim_{d\rightarrow \infty} \mathbb{E}\left[\max_{i \in [n]} \langle \mathbf h, \boldsymbol\theta_i \rangle\right] = f(\rho, \alpha),
\end{equation*}
where $f(\rho, \alpha)$ is the value of the deterministic optimization problem defined in Proposition \ref{prop: utility_approx}. Next, using the scaling of $\lambda_s$, the search cost becomes,
\begin{equation*}
    \lambda_s \log n = \left(\frac{c_s}{d} + o(d^{-1})\right) (\alpha d) = c_s \alpha + o(1).
\end{equation*}
Finally, from Proposition \ref{prop: kl_approx},
\begin{align*}
    D_{\text{KL}} &=  \frac{d-2}{2} \log \left(\frac{1}{1-\rho^2}\right) + \mathcal{E}_{\ref{prop: kl_approx}}(\rho),
\end{align*}
where we use $D_{\text{KL}}$ as a shorthand notation for the KL-divergence. Multiplying by the communication cost parameter $\lambda_c$ and taking the limit,
\begin{align*}
   \lim_{d\rightarrow \infty} \lambda_c D_{\text{KL}} &= \lim_{d\rightarrow \infty} \left(\frac{c_c}{d} + o(d^{-1})\right) \left(\frac{d-2}{2} \log \left(\frac{1}{1-\rho^2}\right) + \mathcal{E}_{\ref{prop: kl_approx}}(\rho)\right) = -\frac{c_c}{2} \log (1-\rho^2),
\end{align*}
where we use the fact that $|\mathcal{E}_{\ref{prop: kl_approx}}(\rho)| \leq \frac{K_{\ref{prop: kl_approx}} \sqrt{d}}{\sqrt{(1-\rho^2)^3}}$, and $K_{\ref{prop: kl_approx}}$ is a universal constant. Combining these terms, the objective function converges to
\begin{equation*}
    \lim_{d \to \infty} \mathcal{P}_d(\kappa, n) = f(\rho, \alpha) - c_s \alpha + \frac{c_c}{2} \log(1-\rho^2).
\end{equation*}
This matches the definition of the objective function in $\mathrm{OPT}_{\mathrm{Joint}}$ (Eq. \eqref{eq:joint_opt}). This shows that the finite-dimensional objective converges uniformly to the asymptotic objective. {Now, under the assumption that the optimizer ($\rho^* , \alpha^*$) lies on a compact set, the maximum also converges, proving the first statement.

As such, to complete the argument, we show that for $c_s > 0$ and $c_c > 0$, any optimizer of $\mathcal{P}_\infty$ must lie in a compact subset of $[0,1) \times [0,\infty)$. Since $f(\rho, \alpha) \leq 1$, the objective satisfies
\begin{align*}
    \mathcal{P}_\infty(\rho, \alpha) \leq 1 - c_s \alpha + \frac{c_c}{2} \log(1-\rho^2).
\end{align*}
Note that $\mathcal{P}_\infty(0, 0) = f(0,0) = 0$ provides a baseline value. For any $(\rho, \alpha)$ to be optimal, we need $\mathcal{P}_\infty(\rho, \alpha) \geq 0$, which requires both $c_s \alpha \leq 1$ and $\frac{c_c}{2} \log(1-\rho^2) > -1$. The first condition gives $\alpha \leq 1/c_s$. The second gives $1-\rho^2 > e^{-2/c_c}$, i.e., $\rho < \sqrt{1 - e^{-2/c_c}} < 1$. Since the optimizer is confined to the compact set $[0, \sqrt{1 - e^{-2/c_c}}] \times [0, 1/c_s]$, the uniform convergence of $\mathcal{P}_d$ to $\mathcal{P}_\infty$ on this set implies convergence of the optima.}

\paragraph{Part 2: Monotonicity}
We denote the asymptotic objective function by $\mathcal{P}_\infty(\rho, \alpha)$ which is given by
\begin{equation*}
    \mathcal{P}_\infty(\rho, \alpha) = f(\rho, \alpha) - c_s \alpha + \frac{c_c}{2} \log(1-\rho^2).
\end{equation*}
For simplicity of the argument, we assume that the optimal solution $(\rho^*, \alpha^*)$ lies in the interior of the feasible region $(0, 1) \times (0, \infty)$. From the first-order necessary conditions for a local maximum, we have
\begin{align}
\label{eq:foc}
    \frac{\partial}{\partial \rho} \mathcal{P}_\infty = \frac{\partial f}{\partial \rho} - \frac{c_c \rho}{1-\rho^2} = 0, && \frac{\partial}{\partial \alpha} \mathcal{P}_\infty = \frac{\partial f}{\partial \alpha} - c_s = 0
\end{align}
For the solution to be a local maximum, the second-order sufficient conditions require that
\begin{align*}
    \frac{\partial^2}{\partial \rho^2} \mathcal{P}_\infty <0, && \frac{\partial^2}{\partial \alpha^2} \mathcal{P}_\infty <0,
\end{align*}
and the Hessian matrix $\mathbf{H}_\infty$ of the objective function $\mathcal{P}_\infty$ with respect to $\rho$ and $\alpha$ be negative definite, where the Hessian is given by
\begin{equation*}
    \mathbf{H}_\infty
    = \begin{pmatrix}
        \frac{\partial^2 f}{\partial \rho^2} - \frac{c_c(1+\rho^2)}{(1-\rho^2)^2} & \frac{\partial^2 f}{\partial \rho \partial \alpha} \\
        \frac{\partial^2 f}{\partial \alpha \partial \rho} & \frac{\partial^2 f}{\partial \alpha^2},
    \end{pmatrix}
\end{equation*}
where second order conditions imply that $|\mathbf{H}_\infty | > 0$. Furthermore, we invoke the property that communication precision and search set size act as substitutes in the production of utility, simply because as $\rho$ increases, the alignment between the agent's recommendation and the user's preference increases. This reduces the necessity of a large search set exponent $\alpha$ to find a high-utility outcome, as such, we have $\frac{\partial^2}{\partial \alpha \partial \rho } \mathcal{P}_\infty = \frac{\partial^2 f}{\partial \rho \partial \alpha} < 0$.

Since the first-order conditions must hold even as $c_c$ changes, we can take the total derivative of both equations in Eq. \eqref{eq:foc} with respect to $c_c$ to get that

\begin{equation*}
    \mathbf{H}_\infty
    \begin{pmatrix}
        \frac{\partial \rho^*}{\partial c_c} \\
        \frac{\partial \alpha^*}{\partial c_c}
    \end{pmatrix}
    +
    \begin{pmatrix}
        \frac{\partial^2}{\partial \rho \partial c_c} \mathcal{P}_\infty \\
        \frac{\partial^2}{\partial \alpha \partial c_c} \mathcal{P}_\infty
    \end{pmatrix}
    =
    \begin{pmatrix}
        0 \\ 0
    \end{pmatrix}.
\end{equation*}
Further, we have that
\begin{align*}
    \frac{\partial^2}{\partial \rho \partial c_c} \mathcal{P}_\infty = \frac{\partial}{\partial c_c} \left( \frac{\partial f}{\partial \rho} - \frac{c_c \rho}{1-\rho^2} \right) = - \frac{\rho}{1-\rho^2}, &&
    \frac{\partial^2}{\partial \alpha \partial c_c} \mathcal{P}_\infty = \frac{\partial}{\partial c_c} \left( \frac{\partial f}{\partial \alpha} - c_s \right) = 0
\end{align*}
Next, we apply Cramer's rule to solve for the linear system, giving us
\begin{equation*}
    \frac{\partial \rho^*}{\partial c_c} = - \frac{1}{|\mathbf{H}_\infty|} \det \begin{pmatrix}
        \frac{-\rho}{1-\rho^2} & \frac{\partial^2}{\partial \alpha \partial \rho } \mathcal{P}_\infty  \\
        0 &
        \frac{\partial^2}{\partial \alpha^2 } \mathcal{P}_\infty
    \end{pmatrix} = \frac{1}{|\mathbf{H}_\infty|} \left( \frac{\rho}{1-\rho^2} \frac{\partial^2}{\partial \alpha^2} \mathcal{P}_\infty \right)< 0,
\end{equation*}
where we use the fact that $|\mathbf{H}_\infty|>0$ and $\frac{\partial^2}{\partial \alpha^2} \mathcal{P}_\infty<0$. Similarly, we get that
\begin{align*}
    \frac{\partial \alpha^*}{\partial c_c} = - \frac{1}{|\mathbf{H}_\infty|} \det \begin{pmatrix}
        \frac{\partial^2}{\partial \rho^2 } \mathcal{P}_\infty & \frac{-\rho}{1-\rho^2} \\
        \frac{\partial^2}{\partial \alpha \partial \rho } \mathcal{P}_\infty & 0
    \end{pmatrix} = \frac{1}{|\mathbf{H}_\infty|} \left( -\frac{\rho}{1-\rho^2} \frac{\partial^2}{\partial \alpha \partial \rho } \mathcal{P}_\infty \right) > 0.
\end{align*}
Next, we repeat the procedure in terms of the search cost parameter $c_s$. We have that
\begin{align*}
    \frac{\partial^2}{\partial \rho \partial c_s} \mathcal{P}_\infty = \frac{\partial}{\partial c_s} \left( \frac{\partial f}{\partial \rho} - \frac{c_c \rho}{1-\rho^2} \right) = 0, &&
    \frac{\partial^2}{\partial \alpha \partial c_s} \mathcal{P}_\infty = \frac{\partial}{\partial c_s} \left( \frac{\partial f}{\partial \alpha} - c_s \right) = -1
\end{align*}
This gives us that
\begin{equation*}
    \frac{\partial \rho^*}{\partial c_s} = \frac{1}{|\mathbf{H}_\infty|} \left( -1 \times \frac{\partial^2}{\partial \alpha \partial \rho } \mathcal{P}_\infty \right)> 0,
\end{equation*}
and also,
\begin{equation*}
    \frac{\partial \alpha^*}{\partial c_s} = \frac{1}{|\mathbf{H}_\infty|} \left( 1 \times \frac{\partial^2}{\partial \rho^2 } \mathcal{P}_\infty \right) < 0.
\end{equation*}
This completes the proof of the monotonicity properties.

\paragraph{Part 3: Search-Only Regime}
To prove this part, we first show that for $c_c > c_s$, the optimal policy for $\mathcal{P}_\infty(\rho,\alpha)$ satisfy that optimal message precision $\rho^* = 0$. From the expression for $f(\rho, \alpha)$ in Proposition \ref{prop: utility_approx}, we know that
\begin{align*}
    f(\rho,\alpha) := \max_{w, x \in (-1,1)} & \ \ \ \rho w + \sqrt{1-\rho^2} \sqrt{1-w^2} x \ \text{ such that } I_{\rho}(w,x) \leq \alpha,
\end{align*}
where $I_{\rho}(w,x)$ is the large deviations rate function, given by
\begin{align*}
    I_{\rho}(w,x) = -\frac{\rho(w-\rho)}{1-\rho^2} - \frac{1}{2} \log(1-w^2)+ \frac{1}{2} \log(1-\rho^2) -\frac{1}{2} \log(1-x^2).
\end{align*}
Using the fact that $-\frac{\rho(w-\rho)}{1-\rho^2} - \frac{1}{2} \log(1-w^2)+ \frac{1}{2} \log(1-\rho^2) \geq 0$ for any $w\in (-1,1)$, we get that
\begin{align*}
    \big\{ (w,x) : |w| < 1, |x|<1, I_{\rho}(w,x) \leq \alpha \big\} \subseteq \Big\{ (w,x) : |w| < 1, |x|<1, -\frac{1}{2} \log(1-x^2)  \leq \alpha \Big\}.
\end{align*}
This implies that
\begin{align*}
    f(\rho,\alpha) \leq \max_{w, x \in (-1,1)} & \ \ \ \rho w + \sqrt{1-\rho^2} \sqrt{1-w^2} x \ \text{ such that } -\frac{1}{2} \log(1-x^2)  \leq \alpha.
\end{align*}
This in turn gives us that,
\begin{align*}
    \mathcal{P}_\infty (\rho, \alpha) &\leq \max_{w, x \in (-1,1)} \ \ \ \rho w + \sqrt{1-\rho^2} \sqrt{1-w^2} x - c_s \alpha - \frac{c_c}{2}\log \frac{1}{1-\rho^2} \ \text{ s.t. } -\frac{1}{2} \log(1-x^2)  \leq \alpha\\
    &= \sqrt{\rho^2 + (1-\rho^2) (1-e^{-2\alpha})} - c_s \alpha - \frac{c_c}{2}\log \frac{1}{1-\rho^2},
\end{align*}
where the second equality is achieved by simply maximizing in terms of $w$ and $x$. It can be shown that when $c_c>c_s$, the $(\rho,\alpha)$ that maximizes the RHS in the above expression satisfy $\rho = 0$. A more similar argument (and in more detail) for the same is presented in Theorem \ref{thm:tilt_opt} and its proof. As such, for $c_c>c_s$, we have
\begin{align*}
    \max_{\rho\in[0,1) , \alpha\geq 0} \mathcal{P}_\infty (\rho, \alpha) \leq \max_{\alpha\geq 0} \sqrt{1-e^{-2\alpha}} - c_s \alpha.
\end{align*}
Further, we also have that,
\begin{align*}
    \max_{\rho\in[0,1) , \alpha\geq 0} \mathcal{P}_\infty (\rho, \alpha) & \geq \max_{ \alpha\geq 0} \mathcal{P}_\infty (0, \alpha) \\
    & = \max_{ \alpha\geq 0} \max_{w, x \in (-1,1)} \Big[\sqrt{1-w^2} x - c_s \alpha \ \text{ s.t. } -\frac{1}{2} \log(1-w^2) -\frac{1}{2} \log(1-x^2)  \leq \alpha \Big]\\
    & = \max_{ \alpha\geq 0}  \max_{x \in (-1,1)} \Big[ x - c_s \alpha \ \text{ s.t. } -\frac{1}{2} \log(1-x^2)  \leq \alpha\Big]\\
    & = \max_{\alpha\geq 0} \sqrt{1-e^{-2\alpha}} - c_s \alpha.
\end{align*}
Combining the above two inequalities, it implies that for $c_c>c_s$,
\begin{align*}
    \max_{\rho\in[0,1) , \alpha\geq 0} \mathcal{P}_\infty (\rho, \alpha) = \max_{\alpha\geq 0} \sqrt{1-e^{-2\alpha}} - c_s \alpha,
\end{align*}
and so $\rho^* = 0$ whenever $c_c>c_s$. Thus, we get that there exists $\bar c_c(c_s)$ such that for $c_c>\bar c_c(c_s)$, we have $\rho^*(c_c,c_s) =0$.
\end{proof}

\end{document}
